% Generated draft: requires mathematical and source review.
\documentclass[11pt]{article}
\usepackage[margin=25mm]{geometry}
\usepackage{fontspec}
\setmainfont{Latin Modern Roman}
\usepackage{amsmath,amssymb,mathtools,unicode-math}
\setmathfont{Latin Modern Math}
\usepackage{enumitem,needspace,xurl,hyperref}
\setlength{\parindent}{0pt}
\setlength{\parskip}{5pt}
\setlength{\emergencystretch}{3em}
\newcommand{\N}{\mathbb{N}}
\newcommand{\Z}{\mathbb{Z}}
\newcommand{\Q}{\mathbb{Q}}
\newcommand{\R}{\mathbb{R}}
\newcommand{\C}{\mathbb{C}}
\newcommand{\F}{\mathbb{F}}
\newcommand{\A}{\mathbb{A}}
\newcommand{\PP}{\mathbb{P}}
\newcommand{\E}{\mathbb{E}}
\newcommand{\eps}{\varepsilon}
\newcommand{\dd}{\,\mathrm{d}}
\begin{document}
\section*{UnsolvedMath ID: 20001964}
\textbf{Source problem code:} catalogue entry
\subsection{Definitions and mathematical statement}
M. Libine's Question 6.7 starts with a noncompact real-algebraic symplectic manifold \((M^{2n},\omega)\), a Hamiltonian action of a compact torus \(T\), a proper semialgebraic moment map \(\mu:M\to\mathfrak t^*\), and compact fixed set \(M^T\). It asks for a geometrically defined equivariant integral satisfying localization and taking values in smooth functions on the whole Lie algebra. The parameter-dependent equivariant form is described as smooth and semialgebraic, not necessarily polynomial; the precise semialgebraic form space, closedness, admissibility and regularization are unspecified.

For a precise conditional boundary formula, use the Cartan convention
\[
 d_X=d-\iota_{X_M},\qquad d_X\alpha(X)=0.
\]
At a regular parameter \(X\) with zero set \(M^T\) and invertible normal Euler factors, choose an invariant metric off the fixed set and set
\[
 \eta_X=\frac{g(X_M,\cdot)}{|X_M|^2},\qquad
 \beta_X=\frac{\eta_X}{d_X\eta_X}
       =-\eta_X\sum_{j\geq0}(\dd\eta_X)^j.
\]
The sum is truncated by form degree and \(d_X\beta_X=1\). For regular moment cutoffs \(M_R=\{\|\mu\|\leq R\}\) containing the fixed set, the aid's Stokes--localization identity has the form
\[
 \int_{M_R}\alpha(X)_{[2n]}
 =\sum_{F\subset M^T}\int_F
       \frac{i_F^*\alpha(X)}{e_T(N_F)(X)}
 +\int_{\partial M_R}(\beta_X\alpha(X))_{[2n-1]}.
\]
Compatible orientations and Euler-class normalizations are required. Thus ordinary cutoff convergence to the ABBV fixed-point expression needs the boundary transgression to vanish; properness and semialgebraicity alone do not ensure this. On compact subsets of a regular parameter chamber, uniform vanishing of this boundary term and all its parameter derivatives suffices for smooth chamberwise convergence. This does not prove smooth extension across singular parameters to the whole Lie algebra.

For a uniformly compactly supported smooth equivariant form, ordinary top-degree integration is a smooth-parameter pushforward and equivariant Stokes holds. For a proposed noncompact cutoff theory, assume the relevant limits exist on an admissible class closed under \(d_X\). It descends to equivariant cohomology exactly when all admissible primitives \(\gamma\) have
\[
 \lim_{R\to\infty}\int_{\partial M_R}\gamma(X)_{[2n-1]}=0,
\]
since \(\int_{M_R}(d_X\gamma)_{[2n]}
=\int_{\partial M_R}\gamma_{[2n-1]}\). This primitive boundary condition must not be presumed from fixed-point compactness.

The explicit obstruction in the aid is the standard circle rotation on \(\C\), with
\[
 \omega=\dd x\wedge\dd y,\quad
 V=-y\partial_x+x\partial_y,\quad
 \mu=(x^2+y^2)/2,\quad
 \lambda=(x\dd y-y\dd x)/2.
\]
For real parameter \(u\), \(d_u=d-u\iota_V\) and
\[
 \alpha(u)=d_u\lambda=\omega-u\mu
\]
is equivariantly exact, polynomial and semialgebraic. The moment map is proper and the fixed set is \(\{0\}\), yet on \(\{\mu\leq R\}\)
\[
 \int_{\{\mu\leq R\}}\alpha(u)_{[2]}=2\pi R,
 \qquad \int_{\{\mu=R\}}\lambda=2\pi R.
\]
The regular-parameter fixed-point contribution is zero because \(\alpha(u)|_0=0\). Thus ordinary cutoffs diverge and cannot give a cohomological pushforward for this entire form class. This does not exclude a differently specified regularized or distributional theory.

The aid reports that later work of Libine gives a distributional partial answer with additional no-zeroes-at-infinity and polynomial-parameter hypotheses, rather than the source's requested smooth function on all of \(\mathfrak t\). The precise theorem and citation are absent.

The extraction also includes Theorem 6.8, attributed to Hausel--Proudfoot (2003): a nondegenerate pairing on \(\widehat H^*_{U(1)}(M)\) given by \(\int_M\alpha\wedge\beta\), described as a form of Poincaré duality for kernel computations. Its hat notation, coefficient ring, integration definition and hypotheses are missing, so it is not asserted here for every preceding noncompact \(M\).

Finally Definition 6.9 calls a hyperkähler manifold hypercompact for \(I\) if it has an \(\omega_I\)-Hamiltonian circle action with proper moment map, finitely many critical points and a minimum, and a compatible complexified action satisfying
\[
 \omega_{\C}=\omega_J+i\omega_K,\qquad
 \lambda^*\omega_{\C}=\lambda\omega_{\C}\quad(\lambda\in\C^*).
\]
The extracted text says critical points, not critical values, and does not explicitly define the needed complexification. These qualifications remain unresolved rather than silently repaired. The following proposed statement is cut off.
\subsection{Short English statement}
A noncompact localization formula has a boundary-at-infinity term. The aid's exact circle-plane example satisfies the geometric assumptions but has a divergent cutoff even for an equivariantly exact form. Compact support or explicit boundary decay can give valid pushforwards; the requested general smooth-valued theory needs more specification.
\subsection{Sources}
\begin{itemize}
\item Permanent UnsolvedMath ID \(20001964\), exact public code catalogue entry; local registry \texttt{aim-geometry-notes.json}, \texttt{notes[301]}, canonical tag question.
\item AIM workshop ``Moment maps and surjectivity in various geometries,'' Question 6.7 attributed to M. Libine, followed by Theorem 6.8 and Definition 6.9: \url{https://aimath.org/WWN/momentmaps/momentmaps.pdf}. No PDF was accessed; the verification note explicitly preserves OCR defects and contextual spillover.
\item Supplied attributions: later distributional localization work of Libine, and Hausel--Proudfoot (2003) for the stated pairing theorem. Full bibliographic details, theorem hypotheses, original literature entry and proposer/year metadata are absent.
\item Research attempt 1, ``Noncompact equivariant integration and the boundary at infinity,'' labelled partial\_result, partially\_solved, proved and low novelty confidence. Its boundary criterion and obstruction are supplied research reports rather than independently checked results.
\end{itemize}
\end{document}
